\chapter{Endliche Halbverbände und irreduzible Elemente}
\label{b45guide:ch:finite}

In diesem Kapitel sind die betrachteten Träger ausdrücklich endlich.
Nichtleere und Null werden dort vorausgesetzt, wo sie benötigt werden.

\section{Von einer Operation zu zwei Operationen}

\begin{theorem}[Eins in einem endlichen Halbverband]
\label{b45guide:thm:finite-top}
Jeder nichtleere endliche Halbverband $(S,\vee)$ besitzt eine Eins,
nämlich
\[
 1=\bigvee S.
\]
Dies setzt keine Null voraus. Dual besitzt jeder nichtleere
endliche infimal gelesene Halbverband ein unterstes Element.
\cite{project}
\end{theorem}

\begin{theorem}[Ergänzung eines endlichen Halbverbandes zum Verband]
\label{b45guide:thm:finite-extension}
Für einen nichtleeren endlichen Halbverband $(S,\vee)$
sind äquivalent:
\begin{enumerate}
\item Es gibt eine Null $0\in S$.
\item Es gibt eine Operation $\wedge$ auf $S$, so daß
$(S,\wedge,\vee)$ ein Verband ist.
\end{enumerate}
Im Existenzfall ist $\wedge$ eindeutig und explizit gegeben durch
\[
 x\wedge y=
 \bigvee\{z\in S\mid z\vee x=x\text{ und }z\vee y=y\}.
\]
Die verknüpfte Menge ist nichtleer, weil sie $0$ enthält.
Das Ergebnis ist ein beschränkter Verband mit Eins $\bigvee S$.
Insbesondere ist jeder nichtleere endliche Verband beschränkt.
\cite{project}
\end{theorem}

\section{Irreduzibilität und eindeutige Nachbarn}

\begin{definition}[Join- und Meet-Irreduzibilität]
\label{b45guide:def:irreducible}
In einem Halbverband $(S,\vee,0)$ heißt $j\in S$
\emph{join-irreduzibel} oder \emph{$\vee$-irreduzibel}, wenn
\[
 j\ne0,\qquad
 \forall u,v\in S\quad
 (j=u\vee v\Rightarrow j=u\text{ oder }j=v).
\]
Die Menge dieser Elemente heißt $J(S)$.

In einem Verband mit Eins heißt $m$ \emph{meet-irreduzibel},
wenn $m\ne1$ und
$m=u\wedge v\Rightarrow m=u$ oder $m=v$ gilt.
Die Menge dieser Elemente heißt $M(L)$.
\end{definition}

\begin{definition}[Unmittelbare Nachbarn]
Für Elemente $a<b$ schreiben wir $a\prec b$, wenn kein Element
streng zwischen $a$ und $b$ liegt. Dann ist $a$ ein
\emph{unterer Nachbar} von $b$ und $b$ ein \emph{oberer Nachbar} von $a$.
\end{definition}

\begin{theorem}[Nachbarkriterium der Irreduzibilität]
\label{b45guide:thm:cover}
Sei $(S,\vee,0)$ endlich und $j\ne0$. Äquivalent sind:
\begin{enumerate}
\item $j$ ist join-irreduzibel.
\item $j$ besitzt genau einen unteren Nachbarn $j_*$.
\item Die Menge $\{x\in S\mid x<j\}$ besitzt ein größtes Element.
\end{enumerate}
Im Existenzfall ist dieses größte Element $j_*$, und
\[
 x<j\quad\Longleftrightarrow\quad x\leq j_* .
\]
Dual ist in einem nichtleeren endlichen Verband $m\ne1$ genau
dann meet-irreduzibel, wenn es einen einzigen oberen Nachbarn
$m^*$ besitzt. Dann gilt $m<x\iff m^*\leq x$.
\cite{project,irreducible}
\end{theorem}

\section{Zerlegung, Trennung und Erzeugung}

\begin{theorem}[Zerlegung in Irreduzible]
\label{b45guide:thm:decomposition}
Sei $(S,\vee,0)$ ein endlicher Halbverband mit Null.
Jedes Element besitzt die Darstellung
\[
 x=\bigvee\{j\in J(S)\mid j\leq x\}.
\]
Für $x=0$ ist die verknüpfte Menge leer und ihr Wert $0$.
Dual gilt für jeden nichtleeren endlichen Verband $L$
\[
 x=\bigwedge\{m\in M(L)\mid x\leq m\},
 \qquad \bigwedge\varnothing=1.
\]
Eine Darstellung durch eine kleinere Auswahl irreduzibler Elemente
muß nicht eindeutig sein. \cite{irreducible}
\end{theorem}

\begin{theorem}[Irreduzible erkennen und trennen Elemente]
\label{b45guide:thm:separation}
Im endlichen Halbverband $(S,\vee,0)$ gilt
\[
 x\leq y
 \quad\Longleftrightarrow\quad
 \forall j\in J(S)\ (j\leq x\Rightarrow j\leq y).
\]
Insbesondere gibt es für $x\not\leq y$ ein $j\in J(S)$
mit $j\leq x$ und $j\not\leq y$.
Die Menge aller unter $x$ liegenden Irreduziblen bestimmt $x$
eindeutig. Ein solcher Halbverband mit mehr als einem Element
besitzt daher mindestens ein join-irreduzibles Element. \cite{project}
\end{theorem}

\begin{theorem}[Kleinste Join-Erzeugermenge]
\label{b45guide:thm:minimal-generators}
Für jeden endlichen Halbverband $(S,\vee,0)$ gilt
\[
 \langle J(S)\rangle_{\vee,0}=S.
\]
Ist $X\subseteq S$ und $\langle X\rangle_{\vee,0}=S$,
so muß $J(S)\subseteq X$ gelten.
Somit ist $J(S)$ die eindeutig kleinste Erzeugermenge
für die Operation $\vee$ unter Mitführung von $0$.
Diese Aussage betrifft die Erzeugung mit einer Operation;
sie behauptet keine kleinste Erzeugermenge unter den
\emph{beiden} Verbandsoperationen gemeinsam.
\end{theorem}

\begin{theorem}[Größenschranken durch Irreduzible]
\label{b45guide:thm:cardinality}
Sei $(S,\vee,0)$ endlich und $r=|J(S)|$. Dann
\[
 r+1\leq |S|\leq2^r.
\]
Insbesondere gilt $J(S)=\varnothing$ genau dann, wenn $S=\{0\}$.
Für einen nichtleeren endlichen Verband gelten außerdem die
dualen Schranken
\[
 |M(L)|+1\leq|L|\leq2^{|M(L)|}.
\]
\end{theorem}

\section{Kanonische Mengendarstellung endlicher Verbände}

\begin{theorem}[Darstellung durch abgeschlossene Irreduziblenmengen]
\label{b45guide:thm:closure-representation}
Sei $L$ ein nichtleerer endlicher Verband und $J=J(L)$. Für
$U\subseteq J$ und $x\in L$ setze
\[
 c(U)=\{j\in J\mid j\leq\bigvee U\},\qquad
 J_x=\{j\in J\mid j\leq x\}.
\]
Dabei gilt $\bigvee\varnothing=0$. Es gelten
\[
 U\subseteq c(U),\qquad
 U\subseteq V\Rightarrow c(U)\subseteq c(V),\qquad
 c(c(U))=c(U).
\]
Die Familie $\mathcal C=\{U\subseteq J\mid c(U)=U\}$ trägt
die Verbandsoperationen
\[
 U\wedge_{\mathcal C}V=U\cap V,\qquad
 U\vee_{\mathcal C}V=c(U\cup V).
\]
Die Abbildung $x\mapsto J_x$ ist ein Isomorphismus $L\to\mathcal C$.
Insbesondere
\[
 J_{x\wedge y}=J_x\cap J_y,\qquad
 J_{x\vee y}=c(J_x\cup J_y).
\]
Die Abschlußbildung $c$ kann in der zweiten Formel im Allgemeinen
nicht entfallen. \cite{nation}
\end{theorem}
